% This is text associated with the open textbook "Open Electromagnetics"
% See immediately below for author, date
% <License info goes here (boilerplate, not read by project compiler)>

%ORB TITLE About This Book
%ORB AUTHOR 0        # S.W. Ellingson, Virginia Tech (ellingson@vt.edu)
%ORB DATE 20191212   # yyyymmdd
%ORB ABSTRACT Traditional preface information (i.e., what it covers, who this book is for, etc)

\label{m0213_Preface-About_This_Book_v2} % <-- Must match name of this file and module directory name.  Don't mess with this!

{\bf Goals for this book.}
This book is intended to serve as a primary textbook for the second semester of a two-semester course in undergraduate engineering electromagnetics.
The presumed textbook for the first semester is \emph{Electromagnetics Vol.~1},\footnote{S.W.~Ellingson, 
\emph{Electromagnetics~Vol.~1}, 
VT~Publishing, 2018. CC BY-SA 4.0. ISBN 9780997920185.\\ 
\url{https://doi.org/10.21061/electromagnetics-vol-1}} 
which addresses the following topics: 
electric and magnetic fields; 
electromagnetic properties of materials;
electromagnetic waves;
and devices that operate according to associated electromagnetic principles including 
resistors, 
capacitors, 
inductors, 
transformers, 
generators,
and
transmission lines.
The book you are now reading -- \emph{Electromagnetics Vol.~2} -- addresses the following topics:
%
\begin{itemize}
\item Chapter~1 (``Preliminary Concepts'') provides a brief summary of conventions for units, notation, and coordinate systems, and a synopsis of electromagnetic field theory from \emph{Vol.~1}.
\item Chapter~2 (``Magnetostatics Redux'') extends the coverage of magnetostatics in \emph{Vol.~1} to include magnetic forces, rudimentary motors, and the Biot-Savart law.
\item Chapter~3 (``Wave Propagation in General Media'') addresses Poynting's theorem, theory of wave propagation in lossy media, and properties of imperfect conductors.
\item Chapter~4 (``Current Flow in Imperfect Conductors'') addresses the frequency-dependent distribution of current in wire conductors and subsequently the AC impedance of wires.  
\item Chapter~5 (``Wave Reflection and Transmission'') addresses scattering of plane waves from planar interfaces.
\item Chapter~6 (``Waveguides'') provides an introduction to waveguide theory via the parallel plate and rectangular waveguides.
\item Chapter~7 (``Transmission Lines Redux'') extends the coverage of transmission lines in \emph{Vol.~1} to include parallel wire lines, the theory of microstrip lines, attenuation, and power-handling capabilities.  The inevitable but hard-to-answer question ``What's so special about $50~\Omega$?'' is addressed at the end of this chapter.
\item Chapter~8 (``Optical Fiber'') provides an introduction to multimode fiber optics, including the concepts of acceptance angle and modal dispersion.
\item Chapter~9 (``Radiation'') provides a derivation of the electromagnetic fields radiated by a current distribution, emphasizing the analysis of line distributions using the Hertzian dipole as a differential element.
\item Chapter~10 (``Antennas'') provides an introduction to antennas, emphasizing equivalent circuit models for transmission and reception and characterization in terms of directivity and pattern.  This chapter concludes with the Friis transmission equation. 
\end{itemize}
%
Appendices covering material properties, mathematical formulas, and physical constants are repeated from \emph{Vol.~1}, with a few additional items.
 
%{\bf What's new:}
%This version of the book is the \emph{second} public release of this book.
%The first release, known as ``Volume~1 (Beta),'' was released in January 2018.
%Improvements from the beta version include the following:
 
{\bf Target audience.}
This book is intended for electrical engineering students in the third year of a bachelor of science degree program.
It is assumed that students have successfully completed one semester of engineering electromagnetics, nominally using \emph{Vol.~1}.
However, the particular topics and sequence of topics in \emph{Vol.~1} are not an essential prerequisite, and in any event this book may be useful as a supplementary reference when a different textbook is used.
It is assumed that readers are familiar with the fundamentals of electric circuits and linear systems, which are normally taught in the second year of the degree program.  
It is also assumed that readers are proficient in basic engineering mathematics, including 
complex numbers,
trigonometry,
vectors,
partial differential equations,
and multivariate calculus.

{\bf Notation, examples, and highlights.}
Section~\ref{m0005_Notation} summarizes the mathematical notation used in this book.
Examples are set apart from the main text as follows:

%%%%%%%%%%%%%%%%%%%%%%%%%%%%%%%%%%%%%%%%%%%%%%%
%ORB EXAMPLE
~\\ \begin{example} 
\begin{mdframed}[backgroundcolor=brown!20,linecolor=brown!20]\raggedright
This is an example.
%\\~\\
%Main body of example...
\end{mdframed}
% note figures will need to go between "\end{mdframed}" and "\end{example}"
\end{example}
%%%%%%%%%%%%%%%%%%%%%%%%%%%%%%%%%%%%%%%%%%%%%%%

``Highlight boxes'' are used to identify key ideas as follows:

%ORB HIGHLIGHT
\begin{mdframed}[backgroundcolor=yellow!20,nobreak=true] 
This is a key idea.
\end{mdframed}

{\bf What are those little numbers in square brackets?}
This book is a product of the 
\hyperref[m0148_Preface-About_OpenEM]{\emph{Open Electromagnetics Project}}. %({\color{red}\url{URL-TBD}}).
This project provides a large number of sections (``modules'') which are assembled (``remixed'') to create new and different versions of the book.
The text ``{\color{blue}\tiny[m0213]}'' that you see at the beginning of this section uniquely identifies the module within the larger set of modules provided by the project.
This identification is provided because different remixes of this book may exist, each consisting of a different subset and arrangement of these modules. 
%Readers can use module identifiers to find additional resources available at the \href{URL-TBD}{Open Electromagnetics website}, which is common to all remixes. 
Prospective authors can use this identification as an aid in creating their own remixes.

{\bf Why do some sections of this book seem to repeat material presented in previous sections?}
In some remixes of this book, authors might choose to eliminate or reorder modules.  
For this reason, the modules are written to ``stand alone'' as much as possible.
As a result, there may be some redundancy between sections that would not be present in a traditional (non-remixable) textbook.
While this may seem awkward to some at first, there are clear benefits: In particular, it never hurts to review relevant past material before tackling a new concept.
And, since the electronic version of this book is being offered at no cost, there is not much gained by eliminating this useful redundancy.

{\bf Why cite Wikipedia pages as additional reading?}
Many modules cite Wikipedia entries as sources of additional information.
Wikipedia represents both the best and worst that the Internet has to offer.
Most educators would agree that citing Wikipedia pages as primary sources is a bad idea, since quality is variable and
content is subject to change over time.
On the other hand, many Wikipedia pages are excellent, and serve as useful sources of relevant information that is not strictly within the scope of the curriculum.  
Furthermore, students benefit from seeing the same material presented differently, in a broader context, and with the additional references available as links from Wikipedia pages. 
We trust instructors and students to realize the potential pitfalls of this type of resource and to be alert for problems.

{\bf Acknowledgments.}
Here's a list of talented and helpful people who contributed to this book:\\
~\\
\emph{The staff of Virginia Tech Publishing, University Libraries, Virginia Tech:}\\
Acquisitions/Developmental Editor \& Project Manager: Anita Walz\\
Advisors: Peter Potter, Corinne Guimont\\
Cover, Print Production: Robert Browder\\
~\\
\emph{Other Virginia Tech contributors:}\\
Accessibility: Christa Miller, Corinne Guimont, Sarah Mease\\
Assessment: Anita Walz\\
~\\
\emph{Virginia Tech Students:}\\
Alt text writer: Michel Comer\\
Figure designers: Kruthika Kikkeri, Sam Lally, Chenhao Wang\\
%Figure designer: Michaela Goldammer\\
%Figure designer: Youmin Qin\\
~\\
\emph{Copyediting:}\\
%Melissa Ashman, Kwantlen Polytechnic University \\
Longleaf Press\\
~\\
\emph{External reviewers:}\\
%Samir El-Ghazaly, University of Arkansas\\
%Stephen Gedney, University of Colorado Denver\\
Randy Haupt, Colorado School of Mines\\
Karl Warnick, Brigham Young University\\
Anonymous faculty member, research university\\
~\\
Also, thanks are due to the students of the Fall 2019 section of ECE3106 at Virginia Tech who used the beta version of this book and provided useful feedback. 

Finally, we acknowledge all those who have contributed their art to Wikimedia Commons (\url{https://commons.wikimedia.org/}) under open licenses, allowing their work to appear as figures in this book.  These contributors are acknowledged in figures and in the ``Image Credits'' section at the end of each chapter. Thanks to each of you for your selfless effort. 
